% Blueprint: pi_1 of the circle is Z
% Created by Numina

\begin{document}

\section{The Fundamental Group of the Circle}

We compute the fundamental group of the complex unit circle
\(\mathbb{S}^1 = \{z \in \mathbb{C} : |z| = 1\}\) (Mathlib's \texttt{Circle}) based at
\(1\), proving it is the infinite cyclic group \(\mathbb{Z}\).

The argument is the classical covering-space computation. The exponential map
\[
  \exp : \mathbb{R} \to \mathbb{S}^1, \qquad \exp(r) = e^{2\pi i r}
\]
(Mathlib's \texttt{Circle.exp}, normalized so that \(\exp(r) = 1\) iff
\(r \in 2\pi\mathbb{Z}\)) is a covering map whose total space \(\mathbb{R}\) is
contractible, hence simply connected. The \emph{monodromy} of this covering --- lifting
a loop based at \(1\) to a path in \(\mathbb{R}\) starting at \(0\) and reading off the
endpoint --- defines the \emph{winding number}, a group isomorphism from the fundamental
group to \(\mathbb{Z}\).

Throughout, \(p = \exp\), the basepoint of \(\mathbb{R}\) is \(0\) (which lies in the fibre
since \(\exp(0) = 1\)), and the fibre over \(1\) is \(p^{-1}(\{1\}) = 2\pi\mathbb{Z}\).
We freely use that \(p\) is a covering map (Mathlib: \texttt{Circle.isCoveringMap\_exp}) and
the monodromy API \texttt{IsCoveringMap.monodromy} together with its compatibility law
\texttt{monodromy\_trans\_apply} and the lift-projection law \texttt{monodromy\_map}. To pass
between an element \(\gamma\) of the fundamental group and its homotopy class of paths we use
the bridge \texttt{FundamentalGroup.toPath} / \texttt{FundamentalGroup.fromPath}; monodromy is
always applied to \(\operatorname{toPath}\gamma\).

\begin{lemma}[Cancelling the period]
    \label{lem:two-pi-cancel}
    \lean{LeanEval.Topology.two_pi_cancel}
    \leanfile{LeanEval/Topology/Pi1Circle.lean}
    \leanok
    For integers \(n, m\), if \(n \cdot 2\pi = m \cdot 2\pi\) as real numbers, then \(n = m\).
\end{lemma}
\begin{proof}
    Since \(2\pi \neq 0\) (\texttt{Real.two\_pi\_pos}, \texttt{Real.pi\_ne\_zero}), right
    cancellation \texttt{mul\_right\_cancel\_0} gives \((n : \mathbb{R}) = (m : \mathbb{R})\), and
    injectivity of the integer cast \texttt{Int.cast\_injective} yields \(n = m\).
\end{proof}

\begin{definition}[Integer parametrization of the fibre]
    \label{def:fiber-equiv-int}
    \lean{LeanEval.Topology.fiberEquivInt}
    \leanfile{LeanEval/Topology/Pi1Circle.lean}
    \uses{lem:two-pi-cancel}
    \leanok
    The fibre \(p^{-1}(\{1\}) = \{r \in \mathbb{R} : \exp(r) = 1\}\) is in bijection with
    \(\mathbb{Z}\) via the map \(\langle r\rangle \mapsto n\), where \(n\) is the unique
    integer with \(r = n \cdot 2\pi\); its inverse is \(n \mapsto \langle n \cdot 2\pi\rangle\).
    Existence and uniqueness of \(n\) follow from the kernel description
    \texttt{Circle.exp\_eq\_one} together with \(2\pi \neq 0\).
\end{definition}

\begin{lemma}[Fibre parametrization on integers]
    \label{lem:fiber-equiv-int-symm}
    \lean{LeanEval.Topology.fiberEquivInt_symm_apply}
    \leanfile{LeanEval/Topology/Pi1Circle.lean}
    \uses{def:fiber-equiv-int}
    \leanok
    For every integer \(n\), the fibre point assigned to \(n\) by the inverse of the parametrization
    of Definition~\ref{def:fiber-equiv-int} has underlying real number \(n\cdot 2\pi\); that is,
    \(\bigl(\mathrm{fiberEquivInt}^{-1}(n) : \mathbb{R}\bigr) = n\cdot(2\pi)\).
\end{lemma}
\begin{proof}
    \uses{def:fiber-equiv-int}
    By construction (Definition~\ref{def:fiber-equiv-int}) the inverse sends \(n\) to the fibre point
    \(\langle n\cdot 2\pi\rangle\), whose underlying real number is \(n\cdot 2\pi\).
\end{proof}

\begin{lemma}[The translated lift is the lift]
    \label{lem:translated-lift-is-lift}
    \lean{LeanEval.Topology.translatedLift_eq_liftPath}
    \leanfile{LeanEval/Topology/Pi1Circle.lean}
    \leanok
    Let \(\gamma_0\) be a path (loop) in \(\mathbb{S}^1\) based at \(1\), let \(a \in \mathbb{R}\)
    satisfy \(\exp(a) = 1\), and let \(e\) be a point of the fibre over \(1\). Writing
    \(\mathrm{liftPath}\,\gamma_0\,e\) (Mathlib's \texttt{IsCoveringMap.liftPath}) for the lift of the
    \emph{path} \(\gamma_0\) starting at \(e\), the translated map
    \(t \mapsto \mathrm{liftPath}\,\gamma_0\,e\,(t) + a\) equals \(\mathrm{liftPath}\,\gamma_0\,(e + a)\),
    the lift of \(\gamma_0\) starting at \(e + a\).
\end{lemma}
\begin{proof}
    The translated map is continuous, sends \(0\) to \(e + a\), and projects to \(\gamma_0\): on the
    one hand \(\exp\circ\mathrm{liftPath}\,\gamma_0\,e = \gamma_0\) by
    \texttt{IsCoveringMap.liftPath\_lifts}, and on the other
    \(\exp(\mathrm{liftPath}\,\gamma_0\,e\,(t) + a) = \exp(\mathrm{liftPath}\,\gamma_0\,e\,(t))\cdot\exp(a)
    = \exp(\mathrm{liftPath}\,\gamma_0\,e\,(t))\) by \texttt{Circle.exp\_add} and \(\exp(a) = 1\). By the
    unique characterization of lifted paths (\texttt{IsCoveringMap.eq\_liftPath\_iff'}) it therefore
    equals \(\mathrm{liftPath}\,\gamma_0\,(e + a)\).
\end{proof}

\begin{lemma}[Translation invariance of monodromy]
    \label{lem:monodromy-translation}
    \lean{LeanEval.Topology.monodromy_translation}
    \leanfile{LeanEval/Topology/Pi1Circle.lean}
    \uses{lem:translated-lift-is-lift}
    \leanok
    Let \(\gamma\) be a homotopy class of loops in \(\mathbb{S}^1\) based at \(1\), let
    \(a \in \mathbb{R}\) satisfy \(\exp(a) = 1\), and let \(e\) be a point of the fibre over
    \(1\) (so \(e + a\) is again in the fibre). Then, as real numbers,
    \[
      \operatorname{monodromy}(\operatorname{toPath}\gamma)(e + a)
        = \operatorname{monodromy}(\operatorname{toPath}\gamma)(e) + a .
    \]
\end{lemma}
\begin{proof}
    \uses{lem:translated-lift-is-lift}
    Destructure the quotient: write \(\operatorname{toPath}\gamma = \llbracket\gamma_0\rrbracket\) for a
    representative path \(\gamma_0\). By definition of monodromy,
    \(\operatorname{monodromy}(\operatorname{toPath}\gamma)(e')\) is the endpoint
    \((\mathrm{liftPath}\,\gamma_0\,e')(1)\) of the lift of \(\gamma_0\) starting at \(e'\). By
    Lemma~\ref{lem:translated-lift-is-lift}, \(\mathrm{liftPath}\,\gamma_0\,(e + a) =
    (t \mapsto \mathrm{liftPath}\,\gamma_0\,e\,(t) + a)\); evaluating at \(t = 1\) gives
    \((\mathrm{liftPath}\,\gamma_0\,(e + a))(1) = (\mathrm{liftPath}\,\gamma_0\,e)(1) + a\), that is,
    \(\operatorname{monodromy}(\operatorname{toPath}\gamma)(e + a) =
    \operatorname{monodromy}(\operatorname{toPath}\gamma)(e) + a\).
\end{proof}

\begin{lemma}[Monodromy on a fibre point, in real coordinates]
    \label{lem:monodromy-val-of-fibre}
    \lean{LeanEval.Topology.monodromy_val_of_fibre}
    \leanfile{LeanEval/Topology/Pi1Circle.lean}
    \uses{lem:monodromy-translation}
    \leanok
    For a class \(\gamma\) in the fundamental group at \(1\) and any fibre point \(e\) over \(1\),
    \[
      \bigl(\operatorname{monodromy}(\operatorname{toPath}\gamma)(e)\bigr).\mathrm{val}
        = \bigl(\operatorname{monodromy}(\operatorname{toPath}\gamma)(\langle 0\rangle)\bigr).\mathrm{val}
          + (e : \mathbb{R}).
    \]
\end{lemma}
\begin{proof}
    \uses{lem:monodromy-translation}
    The fibre point \(e\) equals \(\langle 0\rangle + (e : \mathbb{R})\), since \(\langle 0\rangle\) has
    underlying real number \(0\). Apply Lemma~\ref{lem:monodromy-translation} with base point
    \(\langle 0\rangle\) and \(a = (e : \mathbb{R})\) (which satisfies \(\exp(a) = 1\) because \(e\) lies
    in the fibre over \(1\)), then read off underlying real numbers.
\end{proof}

\begin{definition}[Winding number]
    \label{def:wind}
    \lean{LeanEval.Topology.wind}
    \leanfile{LeanEval/Topology/Pi1Circle.lean}
    \uses{def:fiber-equiv-int}
    \leanok
    For a class \(\gamma\) in the fundamental group of \(\mathbb{S}^1\) at \(1\), the
    \emph{winding number} \(w(\gamma) \in \mathbb{Z}\) is the integer assigned by
    Definition~\ref{def:fiber-equiv-int} to the fibre point
    \(\operatorname{monodromy}(\operatorname{toPath}\gamma)(\langle 0\rangle)\); equivalently,
    \(\operatorname{monodromy}(\operatorname{toPath}\gamma)(\langle 0\rangle) = w(\gamma)\cdot 2\pi\).
    Here \(\langle 0\rangle\) is the chosen lift of the basepoint, valid since \(\exp(0) = 1\)
    (\texttt{Circle.exp\_zero}), and \(\operatorname{toPath}\gamma\) is the path homotopy class of
    \(\gamma\) (\texttt{FundamentalGroup.toPath}).
\end{definition}

\begin{lemma}[Winding number specification]
    \label{lem:wind-spec}
    \lean{LeanEval.Topology.wind_spec}
    \leanfile{LeanEval/Topology/Pi1Circle.lean}
    \uses{def:wind, lem:fiber-equiv-int-symm}
    \leanok
    For a class \(\gamma\) in the fundamental group at \(1\), the underlying real number of the
    fibre point \(\operatorname{monodromy}(\operatorname{toPath}\gamma)(\langle 0\rangle)\) equals
    \(w(\gamma)\cdot 2\pi\).
\end{lemma}
\begin{proof}
    \uses{def:wind, lem:fiber-equiv-int-symm}
    By Definition~\ref{def:wind}, \(\operatorname{monodromy}(\operatorname{toPath}\gamma)(\langle 0\rangle)
      = \mathrm{fiberEquivInt}^{-1}(w(\gamma))\). Taking underlying real numbers and applying
    Lemma~\ref{lem:fiber-equiv-int-symm} with \(n = w(\gamma)\) gives
    \(\bigl(\operatorname{monodromy}(\operatorname{toPath}\gamma)(\langle 0\rangle) : \mathbb{R}\bigr)
      = w(\gamma)\cdot 2\pi\).
\end{proof}

\begin{lemma}[Group law as path concatenation]
    \label{lem:mul-eq-trans}
    \lean{LeanEval.Topology.toPath_mul}
    \leanfile{LeanEval/Topology/Pi1Circle.lean}
    \leanok
    For classes \(\gamma, \delta\) in the fundamental group at \(1\),
    \[
      \operatorname{toPath}(\gamma \cdot \delta)
        = (\operatorname{toPath}\delta).\mathrm{trans}\,(\operatorname{toPath}\gamma).
    \]
\end{lemma}
\begin{proof}
    The fundamental group is the endomorphism monoid \(\operatorname{End}\) of the basepoint in the
    fundamental groupoid, where multiplication is reverse composition,
    \(\gamma \cdot \delta = \delta \mathbin{\gg} \gamma\) (\texttt{CategoryTheory.End.mul\_def}).
    Composition in the fundamental groupoid is path concatenation
    (\texttt{FundamentalGroupoid.comp\_eq}), so
    \(\operatorname{toPath}(\delta \mathbin{\gg} \gamma)
      = (\operatorname{toPath}\delta).\mathrm{trans}\,(\operatorname{toPath}\gamma)\).
\end{proof}

\begin{lemma}[Additivity of winding in real coordinates]
    \label{lem:wind-mul-real}
    \lean{LeanEval.Topology.wind_mul_real}
    \leanfile{LeanEval/Topology/Pi1Circle.lean}
    \uses{def:wind, lem:wind-spec, lem:mul-eq-trans, lem:monodromy-val-of-fibre}
    \leanok
    For all classes \(\gamma, \delta\) in the fundamental group at \(1\),
    \[
      w(\gamma\cdot\delta)\cdot 2\pi = w(\gamma)\cdot 2\pi + w(\delta)\cdot 2\pi .
    \]
\end{lemma}
\begin{proof}
    \uses{lem:wind-spec, lem:mul-eq-trans, lem:monodromy-val-of-fibre}
    By Lemma~\ref{lem:wind-spec} for \(\gamma\cdot\delta\), the left-hand side equals
    \(\bigl(\operatorname{monodromy}(\operatorname{toPath}(\gamma\cdot\delta))(\langle 0\rangle)\bigr).\mathrm{val}\).
    By Lemma~\ref{lem:mul-eq-trans}, \(\operatorname{toPath}(\gamma\cdot\delta)
      = (\operatorname{toPath}\delta).\mathrm{trans}\,(\operatorname{toPath}\gamma)\), so
    \texttt{IsCoveringMap.monodromy\_trans\_apply} (first factor \(\operatorname{toPath}\delta\), second
    \(\operatorname{toPath}\gamma\)) gives
    \[
      \operatorname{monodromy}(\operatorname{toPath}(\gamma\cdot\delta))(\langle 0\rangle)
        = \operatorname{monodromy}(\operatorname{toPath}\gamma)
            \bigl(\operatorname{monodromy}(\operatorname{toPath}\delta)(\langle 0\rangle)\bigr).
    \]
    Applying Lemma~\ref{lem:monodromy-val-of-fibre} with fibre point
    \(e = \operatorname{monodromy}(\operatorname{toPath}\delta)(\langle 0\rangle)\), then
    Lemma~\ref{lem:wind-spec} for \(\gamma\) and for \(\delta\), yields
    \(w(\gamma)\cdot 2\pi + w(\delta)\cdot 2\pi\).
\end{proof}

\begin{lemma}[Winding number is additive]
    \label{lem:wind-mul}
    \lean{LeanEval.Topology.wind_mul}
    \leanfile{LeanEval/Topology/Pi1Circle.lean}
    \uses{lem:wind-mul-real, lem:two-pi-cancel}
    \leanok
    For all classes \(\gamma, \delta\) in the fundamental group at \(1\),
    \(w(\gamma \cdot \delta) = w(\gamma) + w(\delta)\).
\end{lemma}
\begin{proof}
    \uses{lem:wind-mul-real, lem:two-pi-cancel}
    By Lemma~\ref{lem:wind-mul-real},
    \(w(\gamma\cdot\delta)\cdot 2\pi = w(\gamma)\cdot 2\pi + w(\delta)\cdot 2\pi
      = (w(\gamma) + w(\delta))\cdot 2\pi\), so Lemma~\ref{lem:two-pi-cancel} yields
    \(w(\gamma\cdot\delta) = w(\gamma) + w(\delta)\).
\end{proof}

\begin{lemma}[Path classes in \(\mathbb{R}\) are unique]
    \label{lem:quotient-subsingleton}
    \lean{LeanEval.Topology.quotient_subsingleton}
    \leanfile{LeanEval/Topology/Pi1Circle.lean}
    \leanok
    For any \(c \in \mathbb{R}\), the type \(\operatorname{Path.Homotopic.Quotient}(0, c)\) of
    homotopy classes (rel endpoints) of paths in \(\mathbb{R}\) from \(0\) to \(c\) is a
    subsingleton.
\end{lemma}
\begin{proof}
    \(\mathbb{R}\) is a real topological vector space, hence contractible
    (\texttt{RealTopologicalVectorSpace.contractibleSpace}) and therefore simply connected
    (\texttt{SimplyConnectedSpace.ofContractible}). In a simply connected space the quotient of
    paths between two fixed points is a subsingleton
    (\texttt{SimplyConnectedSpace.instSubsingletonQuotient}).
\end{proof}

\begin{definition}[Homotopy class of the lift]
    \label{def:lift-class}
    \lean{LeanEval.Topology.liftClass}
    \leanfile{LeanEval/Topology/Pi1Circle.lean}
    \leanok
    For a class \(\gamma\) in the fundamental group at \(1\), \(\mathrm{liftClass}\,\gamma\) is the
    homotopy class in \(\mathbb{R}\) (rel endpoints) of the lift of \(\operatorname{toPath}\gamma\)
    starting at \(0\); it is a class of paths from \(0\) to
    \(\operatorname{monodromy}(\operatorname{toPath}\gamma)(\langle 0\rangle).\mathrm{val}\).
\end{definition}

\begin{lemma}[Pushing the lift class forward recovers the loop]
    \label{lem:loop-eq-exp-of-lift-class}
    \lean{LeanEval.Topology.map_liftClass}
    \leanfile{LeanEval/Topology/Pi1Circle.lean}
    \leanok
    For a class \(\gamma\) in the fundamental group at \(1\), let \(\mathrm{liftClass}\,\gamma\) denote
    the homotopy class in \(\mathbb{R}\) (rel endpoints) of the lift of \(\operatorname{toPath}\gamma\)
    starting at \(0\). Then the image of \(\mathrm{liftClass}\,\gamma\) under the map induced by \(\exp\)
    equals \(\operatorname{toPath}\gamma\):
    \[
      (\mathrm{liftClass}\,\gamma).\mathrm{map}\,\langle\exp,\_\rangle = \operatorname{toPath}\gamma .
    \]
\end{lemma}
\begin{proof}
    Destructure the quotient: pick a representative path \(\gamma_0\) of \(\operatorname{toPath}\gamma\),
    so \(\mathrm{liftClass}\,\gamma = \llbracket\mathrm{liftPath}\,\gamma_0\,\langle 0\rangle\rrbracket\).
    Since the induced map on classes commutes with the quotient projection
    (\texttt{Path.Homotopic.Quotient.mk\_map}), its image under \(\exp\) is the class of
    \(\exp\circ\mathrm{liftPath}\,\gamma_0\,\langle 0\rangle\), which equals \(\gamma_0\) by
    \texttt{IsCoveringMap.liftPath\_lifts}. Hence the image is
    \(\llbracket\gamma_0\rrbracket = \operatorname{toPath}\gamma\).
\end{proof}

\begin{lemma}[Winding number is injective]
    \label{lem:wind-injective}
    \lean{LeanEval.Topology.wind_injective}
    \leanfile{LeanEval/Topology/Pi1Circle.lean}
    \uses{def:wind, lem:wind-spec, lem:quotient-subsingleton, lem:loop-eq-exp-of-lift-class}
    \leanok
    The winding number map \(w\) is injective.
\end{lemma}
\begin{proof}
    \uses{lem:wind-spec, lem:quotient-subsingleton, lem:loop-eq-exp-of-lift-class}
    Suppose \(w(\gamma) = w(\delta)\). By Lemma~\ref{lem:wind-spec} the lift classes
    \(\mathrm{liftClass}\,\gamma\) and \(\mathrm{liftClass}\,\delta\) share the endpoints \(0\) and
    \(c = w(\gamma)\cdot 2\pi\). By Lemma~\ref{lem:quotient-subsingleton} the type
    \(\operatorname{Path.Homotopic.Quotient}(0, c)\) is a subsingleton, so
    \(\mathrm{liftClass}\,\gamma = \mathrm{liftClass}\,\delta\). Applying the map induced by \(\exp\) and
    Lemma~\ref{lem:loop-eq-exp-of-lift-class} on both sides gives
    \(\operatorname{toPath}\gamma = \operatorname{toPath}\delta\), hence \(\gamma = \delta\) since
    \texttt{FundamentalGroup.toPath} is injective.
\end{proof}

\begin{lemma}[Lift class projects to its loop class]
    \label{lem:exp-comp-topath}
    \lean{LeanEval.Topology.toPath_fromPath_exp}
    \leanfile{LeanEval/Topology/Pi1Circle.lean}
    \leanok
    Let \(L\) be a path in \(\mathbb{R}\) from \(0\) to \(c\) with \(\exp(c) = 1\), so
    \(\exp \circ L\) is a loop based at \(1\) and
    \(\gamma := \operatorname{fromPath}\,\llbracket \exp\circ L\rrbracket\) is an element of the
    fundamental group at \(1\). Then \(\operatorname{toPath}\gamma = \llbracket \exp\circ L\rrbracket\).
\end{lemma}
\begin{proof}
    Immediate from the round trip \(\operatorname{toPath}(\operatorname{fromPath}\,q) = q\)
    (\texttt{FundamentalGroup.toPath}, \texttt{FundamentalGroup.fromPath}) applied to
    \(q = \llbracket \exp\circ L\rrbracket\).
\end{proof}

\begin{lemma}[Monodromy of an exponential loop class]
    \label{lem:exp-comp-monodromy}
    \lean{LeanEval.Topology.monodromy_exp_comp}
    \leanfile{LeanEval/Topology/Pi1Circle.lean}
    \leanok
    Let \(L\) be a path in \(\mathbb{R}\) from \(0\) to \(c\). Then
    \[
      \operatorname{monodromy}(\llbracket \exp\circ L\rrbracket)(\langle 0\rangle) = \langle c\rangle .
    \]
\end{lemma}
\begin{proof}
    The class \(\llbracket \exp\circ L\rrbracket\) is the image of \(\llbracket L\rrbracket\) under the
    map induced by \(\exp\) (\texttt{Path.Homotopic.Quotient.mk\_map}). Since \(L\) starts at \(0\) and
    ends at \(c\), \texttt{IsCoveringMap.monodromy\_map} gives
    \(\operatorname{monodromy}(\llbracket \exp\circ L\rrbracket)(\langle 0\rangle) = \langle c\rangle\).
\end{proof}

\begin{lemma}[Explicit loop and its lift]
    \label{lem:explicit-loop-lift}
    \lean{LeanEval.Topology.explicit_loop_lift}
    \leanfile{LeanEval/Topology/Pi1Circle.lean}
    \uses{lem:exp-comp-topath, lem:exp-comp-monodromy}
    \leanok
    Fix \(n \in \mathbb{Z}\) and let \(L\) be the straight-line path
    \(t \mapsto t \cdot (n\cdot 2\pi)\) in \(\mathbb{R}\) (the \texttt{Path} built from the continuous map
    \(t \mapsto (t : \mathbb{R}) \cdot (n\cdot 2\pi)\) with \(L(0) = 0\), \(L(1) = n\cdot 2\pi\)). Then
    \(\exp(n\cdot 2\pi) = 1\), so \(\gamma_n := \operatorname{fromPath}\,\llbracket \exp\circ L\rrbracket\)
    is a fundamental group element with \(\operatorname{toPath}\gamma_n = \llbracket\exp\circ L\rrbracket\)
    and \(\operatorname{monodromy}(\operatorname{toPath}\gamma_n)(\langle 0\rangle) = \langle n\cdot 2\pi\rangle\).
\end{lemma}
\begin{proof}
    \uses{lem:exp-comp-topath, lem:exp-comp-monodromy}
    By \texttt{Circle.exp\_int\_mul\_two\_pi}, \(\exp(n\cdot 2\pi) = 1\). Lemma~\ref{lem:exp-comp-topath}
    (with \(c = n\cdot 2\pi\)) gives \(\operatorname{toPath}\gamma_n = \llbracket\exp\circ L\rrbracket\),
    and rewriting by it turns Lemma~\ref{lem:exp-comp-monodromy} (with \(c = n\cdot 2\pi\)) into
    \(\operatorname{monodromy}(\operatorname{toPath}\gamma_n)(\langle 0\rangle) = \langle n\cdot 2\pi\rangle\).
\end{proof}

\begin{lemma}[Winding number is surjective]
    \label{lem:wind-surjective}
    \lean{LeanEval.Topology.wind_surjective}
    \leanfile{LeanEval/Topology/Pi1Circle.lean}
    \uses{def:wind, lem:explicit-loop-lift, lem:wind-spec, lem:two-pi-cancel}
    \leanok
    The winding number map \(w\) is surjective.
\end{lemma}
\begin{proof}
    \uses{lem:explicit-loop-lift, lem:wind-spec, lem:two-pi-cancel}
    Fix \(n \in \mathbb{Z}\) and take \(\gamma_n\) from Lemma~\ref{lem:explicit-loop-lift}, so
    \(\operatorname{monodromy}(\operatorname{toPath}\gamma_n)(\langle 0\rangle) = \langle n\cdot 2\pi\rangle\).
    By Lemma~\ref{lem:wind-spec} this fibre point also has underlying real number
    \(w(\gamma_n)\cdot 2\pi\), so \(w(\gamma_n)\cdot 2\pi = n\cdot 2\pi\), and
    Lemma~\ref{lem:two-pi-cancel} gives \(w(\gamma_n) = n\). Hence \(w\) is surjective.
\end{proof}

\begin{definition}[Winding-number isomorphism]
    \label{def:wind-equiv}
    \lean{LeanEval.Topology.windEquiv}
    \leanfile{LeanEval/Topology/Pi1Circle.lean}
    \uses{def:wind, lem:wind-mul, lem:wind-injective, lem:wind-surjective}
    \leanok
    The winding number, regarded as a map to \(\operatorname{Multiplicative}\mathbb{Z}\) via
    \(\gamma \mapsto \operatorname{ofAdd}(w(\gamma))\), is a group isomorphism
    \[
      \text{(fundamental group of } \mathbb{S}^1 \text{ at } 1) \;\cong\; \operatorname{Multiplicative}\mathbb{Z}.
    \]
    It is a homomorphism by Lemma~\ref{lem:wind-mul} and bijective by
    Lemmas~\ref{lem:wind-injective} and~\ref{lem:wind-surjective}; one packages it with
    \texttt{MulEquiv.ofBijective}.
\end{definition}

\begin{theorem}[Fundamental group of the circle]
    \label{thm:pi1-circle}
    \lean{LeanEval.Topology.pi1_circle_mulEquiv_int}
    \leanfile{LeanEval/Topology/HomotopyGroups.lean}
    \uses{def:wind-equiv}
    \leanok
    The first homotopy group of the complex unit circle \(\mathbb{S}^1\) based at \(1\) is
    isomorphic, as a group, to \(\operatorname{Multiplicative}\mathbb{Z}\):
    \[
      \pi_1(\mathbb{S}^1, 1) \;\cong\; \operatorname{Multiplicative}\mathbb{Z}.
    \]
\end{theorem}
\begin{proof}
    \uses{def:wind-equiv}
    Mathlib's \texttt{HomotopyGroup.pi1MulEquivFundamentalGroup} gives a group isomorphism
    between the first homotopy group \(\pi_1(\mathbb{S}^1, 1)\) and the fundamental group at
    \(1\). Composing it with the winding-number isomorphism of
    Definition~\ref{def:wind-equiv} produces the desired isomorphism, whence the carrier is
    nonempty.
\end{proof}

\end{document}
